\fullpageexercises[Exponential Maps Across Contexts]
{
\section*{\normalsize How does each general definition specialize to \(\exp(z) = e^z\)?}

\begin{enumerate}[leftmargin=*,itemsep=2pt]

    \item \textbf{Formal Power Series.}
    The series \( \exp(x) = \sum_{n=0}^\infty x^n/n! \) in any commutative \( \mathbb{Q} \)-algebra, evaluated over \( \mathbb{C} \), converges to \( e^z \colon \mathbb{C} \to \mathbb{C}^* \) with \( \exp(z+w) = \exp(z)\exp(w) \).

    \item \textbf{Lie Theory.}
    The exponential maps \( X \in \mathfrak{g} \) to the time-1 flow of its left-invariant vector field on \( G \). For \( G = \mathbb{C}^* \), \( \mathfrak{g} \cong \mathbb{C} \), the flow of \( V_X(z) = Xz\,\tfrac{d}{dz} \) is \( z \mapsto e^{Xt}z \), so \( \exp(X) = e^X \).

    \item \textbf{Eigenfunction of a Derivation.}
    Solve \( Df = \lambda f \) with \( f(0)=1 \). For \( D = d/dz \) and \( \lambda=1 \), the unique solution is \( f(z) = e^z \).

    \item \textbf{Microlocal Analysis.}
    Applying a differential operator \( D \) to the probe \( e^{i\xi \cdot x} \) extracts its symbol: \( D(e^{i\xi \cdot x}) = P(x,\xi)\,e^{i\xi \cdot x} \). On \( \mathbb{C} \), with \( D = d/dz \) and the probe \( e^{z} \), the symbol is simply \( 1 \)—the exponential is both eigenfunction and probe.

    \item \textbf{Sheaf Theory.}
    In the exact sequence \( 0 \to 2\pi i\,\mathbb{Z} \to \mathcal{O}_M \to \mathcal{O}_M^* \to 0 \), the map \( f \mapsto e^f \) sends additive to multiplicative sheaves. On \( M = \mathbb{C} \): \( e^z \colon \mathbb{C} \to \mathbb{C}^* \), with kernel \( 2\pi i\mathbb{Z} \) obstructing a global logarithm.

    \item \textbf{Algebraic Homomorphism.}
    A continuous group homomorphism \( \phi \colon (\mathbb{C},+) \to (\mathbb{C}^*,\times) \) satisfying \( \phi(x+y) = \phi(x)\phi(y) \) must take the form \( \phi(z) = e^{\lambda z} \). Setting \( \lambda = 1 \) gives the classical exponential.

    \item \textbf{Category Theory.}
    The exponential object \( Y^X \) in a CCC satisfies \( \mathrm{Hom}(A \times X,Y) \cong \mathrm{Hom}(A,Y^X) \). With \( X = \mathbb{C} \), \( Y = \mathbb{C}^* \): the map \( (a,x) \mapsto e^{ax} \) corresponds to \( a \mapsto (x \mapsto e^{ax}) \), placing \( e^z \) in \( (\mathbb{C}^*)^{\mathbb{C}} \).

    \item \textbf{Riemannian Geometry.}
    The exponential sends \( v \in T_pM \) to the geodesic endpoint \( \gamma_v(1) \). On \( \mathbb{C}^* \) with the left-invariant metric \( \langle u,v\rangle_z = \mathrm{Re}(u\bar{v})/|z|^2 \), geodesics through \( 1 \) are \( t \mapsto e^{tX} \), giving \( \exp_1(X) = e^X \).

\end{enumerate}
}
